\section{Conclusion}
\label{sec:conclusion}

Providing dual-runtime execution for edge intrusion detection requires unequivocal guarantees of runtime equivalence. Our findings establish that standard traffic replay is dangerously insufficient: an unverified float64 traversal that flips 328 labels under our boundary-directed probe protocol slips through conventional $2{,}000$-row replay gates $93.5\%$ of the time. 

By synthesizing input vectors that climb tree ancestor paths into feasible bounding hyper-rectangles, our \emph{Boundary-Directed Parity Testing} protocol (Algorithm~\ref{alg:boundary_probe}) provides a deterministic, zero-dataset release gate that exposes rounding vulnerabilities up to $34\times$ more densely than traffic replay. Furthermore, our cross-paradigm analysis reveals that modern histogram binning inherently concentrates thresholds into dense shared bins, amplifying floating-point rounding errors into cascading label flips that affect up to $1.58\%$ of boundary decisions in XGBoost and LightGBM models. In contrast, our float32 lockstep traversal provably guarantees exact, bitwise parity ($0$ flips) across all evaluated paradigms and datasets.

Complemented by \emph{Cache-Blocked Lockstep Traversal with Exact Early Exit} (Algorithm~\ref{alg:lockstep_early_exit}), our in-process NumPy engine achieves an $80\times$ speedup at unit batch size on physical NVIDIA Jetson Orin Nano boards, delivering cold-start verdicts in only $357\,\text{ms}$ while reducing tree evaluation depth by $92.9\%$ on suspicious traffic. Finally, by decoupling in-process serving throughput ($21{,}191\,\text{flows/s}$) from network ingestion constraints, this work provides a rigorous theoretical foundation and practical systems blueprint for verifiable, high-throughput intrusion detection at the network edge.
